% Imporditud: tools/import_eksperimendid.py
% Allikas: Eesti füüsikaolümpiaadi eksperimendi ülesannete
%   kogu (Rael Kalda, 2023), ülesanne Ü74, lahendus L74.
\setAuthor{Erkki Tempel}
\setRound{piirkonnavoor}
\setYear{2021}
\setNumber{G E1}
\setDifficulty{2}
\setTopic{dünaamika}
\setVahendid{kolm A4 paberilehte, joonlaud}
\setPunktid{10}
\setAllikas{Eksperimendi kogu Ü74, lahendus lk 66}
\setLahendusseis{toores_ilma_skeemita}

\prob{Paber}
Leidke hõõrdetegur kahe paberilehe vahel. Laua liigutamine ei ole lubatud.

\hint

\solu
Hõõrdeteguri saame määrata kasutades kaldpinna meetodit. Voldime ühe/kaks paberit kokku, et tekitada paksem paber, mida saab kasutada kaldpinnana. Kaldpinna peale asetame kolmanda paberi (kokkuvoldutuna) ning leiame minimaalse nurga, mille korral paber hakkab kaldpinnal libisema. Hõõrdeteguri kahe paberilehe vahel saame leida seosest $\mu$ = tan $\alpha$.

\textit{Žürii hindamisskeem on kogumikus lk 66; siia ei ole seda ümber kirjutatud.}
\probend
